\chapter{Clones: Fazendo Cópias}

\begin{objetivos}
\begin{itemize}
\item Entender o que são clones;
\item Criar e apagar clones com blocos de Controle;
\item Fazer uma chuva de objetos usando um só ator.
\end{itemize}
\end{objetivos}

\section{O que é um clone?}

Imagine ter dez maçãs caindo do céu. Em vez de criar dez atores, você cria \textbf{um só} e manda o Scratch fazer \textbf{cópias} dele. Cada cópia é um \textbf{clone}.

\begin{pensaassim}[Pense assim: a máquina de xerox]
Quando você tira uma cópia de um desenho, ganha uma folha igual à original. Depois, pode riscar e pintar a cópia sem estragar o desenho de verdade. Clone é isso: uma cópia do ator que sai com o mesmo código e pode ir para outro lugar, sem atrapalhar o original.
\end{pensaassim}

\begin{vocabulario}[Vocabulário: clone]
\textbf{Clone} é uma cópia temporária de um ator. Ele nasce com o mesmo código e a mesma aparência do original, mas passa a viver a sua própria vida. Quando o jogo para, todos os clones desaparecem.
\end{vocabulario}

\section{Os três blocos dos clones}

\begin{center}
\renewcommand{\arraystretch}{1.8}
\begin{tabular}{@{}ll@{}}
\bloco[scControl]{crie um clone de \menu{mim mesmo}} & faz uma cópia \\
\bloco[scControl]{quando eu começar como um clone} & código que só os clones obedecem \\
\bloco[scControl]{apague este clone} & o clone some para sempre \\
\end{tabular}
\end{center}

\begin{dica}
Pense assim: o \textbf{original} é o criador de clones e os \textbf{clones} são os trabalhadores. O original fica no topo da tela, só criando; os clones é que se mexem.
\end{dica}

\begin{missao}[Missão 15: Chuva de maçãs]
Faça maçãs caírem de lugares diferentes do topo do palco, a cada meio segundo, usando clones.
\end{missao}

Primeiro, o \textbf{original} fica escondido e cria clones sem parar:

\begin{codigo}[Ator original: a fábrica de clones]
\begin{pilha}
\bk[scEvents]{quando \bandeira\ for clicado}
\bk[scLooks]{esconda}
\sbc[scControl]{sempre}
\bk[scMotion]{vá para x: \reporter{número aleatório entre \entrada{$-200$} e \entrada{200}} y: \entrada{170}}
\bk[scControl]{crie um clone de \menu{mim mesmo}}
\bk[scControl]{espere \entrada{0.5} segundos}
\sbcend
\end{pilha}
\end{codigo}

Cada \textbf{clone} aparece, cai até o chão ($y$ menor que $-160$) e some:

\begin{codigo}[Clone: a maçã que cai]
\begin{pilha}
\bk[scControl]{quando eu começar como um clone}
\bk[scLooks]{mostre}
\sbc[scControl]{repita até que \booleano{\variavel{y} $<$ \entrada{$-160$}}}
\bk[scMotion]{mude y em \entrada{$-5$}}
\sbcend
\bk[scControl]{apague este clone}
\end{pilha}
\end{codigo}

\begin{cuidado}
Sempre dê um jeito de \textbf{apagar} os clones! Se eles se acumularem sem parar, o projeto fica pesado e o Scratch trava (o limite é de 300 clones ao mesmo tempo).
\end{cuidado}

\begin{curiosidade}
Muitos jogos famosos usam clones para tiros, inimigos e moedas. Uma única ``bala'' desenhada vira centenas de tiros na tela!
\end{curiosidade}

\section{Hora de praticar}

\begin{exercicio}[Original e clone]
Faça o gatinho, ao ser clicado, criar um clone que gira 90 graus e anda 50 passos.
\end{exercicio}

\begin{solucao}
\begin{pilha}
\bk[scEvents]{quando este ator for clicado}
\bk[scControl]{crie um clone de \menu{mim mesmo}}
\end{pilha}
\begin{pilha}
\bk[scControl]{quando eu começar como um clone}
\bk[scMotion]{gire \entrada{90} graus}
\bk[scMotion]{mova \entrada{50} passos}
\end{pilha}
\end{solucao}

\begin{exercicio}[Clone que some]
Faça o clone do exercício anterior desaparecer (ser apagado) 3 segundos depois de nascer.
\end{exercicio}

\begin{solucao}
\begin{pilha}
\bk[scControl]{quando eu começar como um clone}
\bk[scMotion]{gire \entrada{90} graus}
\bk[scMotion]{mova \entrada{50} passos}
\bk[scControl]{espere \entrada{3} segundos}
\bk[scControl]{apague este clone}
\end{pilha}
\end{solucao}

\begin{exercicio}[Cinco estrelas ao acaso]
Crie 5 clones de um ator, cada um em uma posição aleatória do palco.
\end{exercicio}

\begin{solucao}
\begin{pilha}
\bk[scEvents]{quando \bandeira\ for clicado}
\sbc[scControl]{repita \entrada{5} vezes}
\bk[scMotion]{vá para x: \reporter{número aleatório entre \entrada{$-200$} e \entrada{200}} y: \reporter{número aleatório entre \entrada{$-150$} e \entrada{150}}}
\bk[scControl]{crie um clone de \menu{mim mesmo}}
\sbcend
\end{pilha}
\end{solucao}

\begin{desafio}[Pegando as maçãs]
Combine a Missão 15 com uma variável \textbf{pontos}. Crie um ator \textbf{cesta} que anda com as setas. Quando uma maçã (clone) tocar na cesta, some e dá 1 ponto.
\end{desafio}

\begin{solucao}
No clone da maçã, acrescente dentro do laço de queda:
\begin{pilha}
\bk[scEvents]{quando \bandeira\ for clicado}
\bk[scVariables]{defina \menu{pontos} para \entrada{0}}
\end{pilha}
\begin{pilha}
\bk[scControl]{quando eu começar como um clone}
\bk[scLooks]{mostre}
\sbc[scControl]{repita até que \booleano{\variavel{y} $<$ \entrada{$-160$}}}
\bk[scMotion]{mude y em \entrada{$-5$}}
\sbc[scControl]{se \booleano{tocando em \menu{cesta}?} então}
\bk[scVariables]{mude \menu{pontos} em \entrada{1}}
\bk[scControl]{apague este clone}
\sbcend
\sbcend
\bk[scControl]{apague este clone}
\end{pilha}
\end{solucao}
